% § 4 — Materials: Corpus, Thematic Layer and Clause Templates
\section{Materials: Corpus, Thematic Layer and Clause Templates}
\label{sec:materials}
\label{sec:layer}
\label{sec:templates}
\textbf{Corpus.} \corpus{} comprises 50 online terms of service split into 9,414 sentences, of which 1,032 (11.0\,\%) carry at least one of 1,137 unfairness labels. The 17 held-out contracts contain 4,078 sentences, 440 of them labelled unfair.

\textbf{Thematic layer (companion short paper).} The clause themes come from the resource built in our companion short paper~\cite{elazhar2026thematic}, which we reuse here unchanged. Every one of the 9,414 sentences carries one primary theme and zero or more secondary themes from a closed vocabulary of 20 legal themes: the four themes where unfairness concentrates (governing law, modification of terms, termination, limitation of liability); the contractual machinery (preamble and scope, eligibility and account, fees and payment, warranty disclaimer, arbitration and disputes, third-party services, privacy and data, DMCA notices); the content and conduct themes (licence and IP, user content, acceptable use, feedback, promotions); and the framing themes (communications, meta, miscellaneous boilerplate). Three trained annotators labelled each sentence independently from a written codebook; four language-model judges labelled the same sentences but were kept out of the reference, which is resolved strictly between annotators by a cascade (strict agreement 46.8\,\%, majority 48.3\,\%, human arbitration 4.9\,\%), with individual votes preserved~\cite{krippendorff2018content,passonneau2006masi,artstein2008intercoder,braun2024differ}. Agreement is $\alpha$-MASI 0.658 on the 20 themes and 0.725 on their 11-theme consolidation, whose merges follow observed confusions but never cross unfairness strata; the leave-one-annotator-out human ceiling is $\kappa$ 0.859. The short paper measures the reliability of this layer; the present paper uses it three times: the consensus primary theme delimits clauses (maximal runs of sentences sharing a theme) and selects the action inventory the extractor sees; each query is scoped to the theme under which its item is found (Table~\ref{tab:greylist}); and the theme alone is the first baseline that any structured method must beat. Themes are not unfairness labels, but they are far from independent of them: Table~\ref{tab:themes} gives, for each theme, the share of sentences carrying a \corpus{} label. The four high-stratum themes carry lifts between 3.4 and 5.0 over the 11.0\,\% base rate; the framing themes are almost never unfair.

\begin{table}[t]
\caption{Unfairness by clause theme on the 50 \corpus{} contracts (consensus theme; $n$ sentences, share carrying a \corpus{} label, lift over the 11.0\,\% base rate). Themes with fewer than 20 sentences are omitted.}
\label{tab:themes}
\centering
\footnotesize\setlength{\tabcolsep}{3.5pt}%
\begin{tabular}{@{}lrrr@{\quad}lrrr@{}}
\toprule
Theme & $n$ & unfair & lift & Theme & $n$ & unfair & lift \\
\midrule
Governing law & 164 & 54.3\,\% & 4.95 & Fees and payment & 1,077 & 5.5\,\% & 0.50 \\
Modification of terms & 392 & 49.0\,\% & 4.47 & Warranty disclaimer & 613 & 5.4\,\% & 0.49 \\
Termination & 422 & 43.4\,\% & 3.96 & Eligibility and account & 761 & 4.2\,\% & 0.38 \\
Limitation of liability & 604 & 37.2\,\% & 3.40 & Privacy and data & 278 & 4.0\,\% & 0.36 \\
Arbitration and disputes & 847 & 10.5\,\% & 0.96 & Licence and IP & 1,225 & 3.4\,\% & 0.31 \\
User content & 658 & 9.7\,\% & 0.89 & Acceptable use & 911 & 3.4\,\% & 0.31 \\
Preamble and scope & 677 & 9.5\,\% & 0.86 & Miscellaneous boilerplate & 586 & 0.9\,\% & 0.08 \\
Third-party services & 464 & 6.5\,\% & 0.59 & Meta, communications & 184 & 0.0\,\% & 0.00 \\
\bottomrule
\end{tabular}
\end{table}

\textbf{Template extraction.} Templates are extracted automatically; clauses longer than eight sentences are split into windows of at most eight. A language model (Claude Opus~5) receives the clause sentences, the clause theme, the closed action inventory for that theme and the output schema, and returns zero or more templates, each citing the sentences that support it. It runs in an isolated context that contains neither the queries nor any reference label. Outputs are checked for schema conformance and structural validity (cited sentences exist, actions belong to the inventory); a lexical anchoring check then resets to \emph{not stated} any condition, notice or remedy value that no cue in the cited sentences supports. Templates were not reviewed by legal experts in this study (Section~\ref{sec:discussion}). On a development pilot of 100 clauses extracted three times, 98\,\% of outputs conformed to the schema, the decisive fields (actor, modality, action, condition, notice, remedy) agreed across runs with a mean Jaccard index of 0.81, and the sentences flagged by the frozen queries agreed across runs with Fleiss' $\kappa = 0.86$. The held-out contracts were extracted twice. The pre-registered analysis uses the first pass: 1,087 clause windows, all schema-conformant, 99.9\,\% structurally and semantically valid, yielding 2,809 templates (2.58 per clause), with 9.7\,\% of windows carrying none. The second pass, evaluated with the same frozen analysis, serves as a sensitivity check: its decisive fields agree with those of the first with a mean Jaccard index of 0.77, and the sentences flagged by the queries agree with Fleiss' $\kappa = 0.82$.
